\section*{الفصل \ref{ch:g1:light-and-shadows} --- الضوء والظلال}

\begin{solution}{exo:g1:light-and-shadows:1}
المنابع: لهب الشمعة، واليراعة، والمصباح اليدوي. وغير المنابع:
القمر، والمرآة، والجدار الأبيض --- فهي لا تفعل إلا أن تردّ \omterm{def:g1:light-and-shadows:source}{الضوء}
الواقع عليها.
\end{solution}

\begin{solution}{exo:g1:light-and-shadows:2}
\omterm{def:g1:light-and-shadows:source}{الضوء} هو الناقص. فليس في الغرفة منبع، ولذلك لا يقع \omterm{def:g1:light-and-shadows:source}{ضوء} على
الأشياء، والأشياء التي لا \omterm{def:g1:light-and-shadows:source}{ضوء} عليها لا تُرى.
\end{solution}

\begin{solution}{exo:g1:light-and-shadows:3}
\omterm{def:g1:light-and-shadows:source}{منبع ضوء}، وحاجب، وجدار (أو حاجز أو أرض) --- بهذا الترتيب على
استقامة واحدة: المنبع يضيء، والحاجب يحجب، والجدار يُظهر الرقعة
المعتمة.
\end{solution}

\begin{solution}{exo:g1:light-and-shadows:4}
\omterm{def:g1:light-and-shadows:shadow}{الظلّ} عن يمينك --- في الجهة المقابلة للشمس. وبعد أن تستدير
وتواجه الجهة المعاكسة تصير الشمس عن يمينك، فيقع \omterm{def:g1:light-and-shadows:shadow}{الظلّ} عن يسارك:
دائمًا في الجهة البعيدة عن الشمس.
\end{solution}

\begin{solution}{exo:g1:light-and-shadows:5}
تحريك اللعبة \emph{نحو المصباح} يجعل \omterm{def:g1:light-and-shadows:shadow}{الظلّ} هائلًا؛ وتحريكها
\emph{نحو الجدار} يعيد \omterm{def:g1:light-and-shadows:shadow}{الظلّ} إلى حجم اللعبة.
\end{solution}

\begin{solution}{exo:g1:light-and-shadows:6}
لا. \omterm{def:g1:light-and-shadows:shadow}{الظلّ} ليس شيئًا مصنوعًا من مادة --- إنما هو موضع غاب عنه
\omterm{def:g1:light-and-shadows:source}{الضوء} لأن حاجبًا يمنع \omterm{def:g1:light-and-shadows:source}{الضوء} عنه. وحيث لا شيء، لا شيء يُلتقط.
\end{solution}

\begin{solution}{exo:g1:light-and-shadows:7}
الشمس منبع أقوى بكثير من أي مصباح: فضوءها من القوة بحيث يؤذي
العينين أذًى دائمًا، لا أن يبهرهما لحظة فحسب.
\end{solution}

\begin{solution}{exo:g1:light-and-shadows:8}
في آخر العصر، حين تنخفض الشمس. وينبغي أن يُظهر الرسم ظلًّا
قصيرًا عند أصل الشجرة عند الظهر، وظلًّا طويلًا ممتدًّا بعد
الظهر، وكلاهما مبتعد عن الشمس.
\end{solution}

\begin{solution}{exo:g1:light-and-shadows:9}
سام هو المصيب. فقد بقي \omterm{def:g1:light-and-shadows:shadow}{الظلّ} مقابلًا للشمس تمامًا طوال الوقت.
وإنما استدارت ميا، فتغيّر موضع \omterm{def:g1:light-and-shadows:shadow}{الظلّ} \emph{بالقياس إلى أمامها
وخلفها} --- أما على الأرض فلم يتحرّك البتة.
\end{solution}
